\documentclass[12pt]{article}
\usepackage[utf8]{inputenc}
\usepackage[T1]{fontenc}
\usepackage{hyperref}
\title{Tamang, Eastern Gorkha Language Resources for K-12 Curriculum Development}
\author{Compiled by Claude Sonnet 5 (claude-sonnet-5) for Andr\'as Kornai}
\date{}
\begin{document}
\maketitle

\section{Introduction}
Eastern Gorkha Tamang (ISO 639-3 tge, Glottocode east2346) is a Sino-Tibetan language of Nepal in the Tamangic branch. Asia Harvest notes that speakers may refer to themselves as Gurung. Little variety-specific material was found.

\section{Primary Educational Resources}
None found specific to this variety.

\section{Linguistic and Reference Materials}
\begin{itemize}
\item Glottolog entry: \url{https://glottolog.org/resource/languoid/id/east2346.bigmap.html}
\item Tamang language (Wikipedia, covers all Tamang varieties): \url{https://en.wikipedia.org/wiki/Tamang_language}
\item Asia Harvest people-group profile (lexical similarity to other Tamang varieties): \url{https://people-groups.asiaharvest.org/Nepal/TamangEasternGorkha.pdf}
\item Scribd-hosted Tamang sociolinguistic survey of Gorkha district, not variety-specific and of unverified provenance: \url{https://www.scribd.com/document/736854705/Chapter}
\end{itemize}

\section{Religious and Community Resources}
No scripture or community organisation specific to this variety was found.

\section{Digital and Audio Resources}
Global Recordings Network has a language page with audio recordings: \url{https://globalrecordings.net/en/language/tge}

\section{Curriculum Development Context}
No dictionary, grammar or school programme for this variety was found. Materials for other Tamang varieties may be adaptable, but the lexical similarity figures (about 63--79 percent) indicate that adaptation would need checking with speakers.

\end{document}